%  GMN2 §10 SU(2) 规范理论：$N_f=0$、$N_f=1$、$N_f=2$ 的第一种实现（原文行 7108–7477）
\section{$SU(2)$ 规范理论}\label{section:SU2Examples}

我们在第 9 节% TODO 合稿改 \ref{section:Superconformal}
中刚刚研究过的那些例子只有 BPS 超多重态，因此当 $\vartheta$ 变化时，WKB 三角剖分
$T_\WKB(\vartheta)$ 只呈现出简单的 flip。接下来我们想考察一些带有 BPS 矢量多重态的例子，
以及与之相应的 \emph{juggle} 变换。我们将研究一组众所周知的理论：带 $N_f=0, 1, 2, 3, 4$
的 $SU(2)$ 规范理论。它们谱的某些方面已经被很好地理解。谱在各种参数区域中的显式结果，
特别可参见 \cite{Bilal:1996sk,Ferrari:1996sv,Ferrari:1997gu}；用 $F$-理论进路的，
可参见 \cite{Banks:1996nj,Douglas:1996js,Gaberdiel:1998mv,Mikhailov:1998bx,DeWolfe:1998bi}。
然而，关于它们的 BPS 谱的完整描述尚不存在。本节将朝着这样一种完整描述取得进展，
给出参数空间许多区域中的谱。事实上，我们的方法可以用来得到模空间任意一点处的谱。

\subsection{$N_f=0$}\label{sec:nf0}

这是 Seiberg--Witten 理论的典型例子。它可以如第 3 节% TODO 合稿改 \ref{sec:physics}
那样，用两根张在两片 NS5-膜之间的 D4-膜构造出来（见方程 (3.130)% TODO 合稿改 \eqref{eq:Standard-SU2-curve}
及其后）。二次微分简单地就是
\begin{equation}
\lambda^2 = \left( \frac{\Lambda^2}{z^3} + \frac{2u}{z^2} + \frac{\Lambda^2}{z} \right) dz^2.
\end{equation}
相应的 Hitchin 系统在 $z=0, \infty$ 处具有可能的最温和的不规则奇点，每个奇点只射出一条
WKB 射线（对应于第 8 节% TODO 合稿改 \ref{sec:irregular}
记号中的 $L=1$）。共有两个转向点，$z_{tp}^\pm = -\frac{u}{\Lambda^2} \pm \sqrt{(\frac{u}{\Lambda^2})^2-1}$，
它们在 $u = \pm \Lambda^2$ 时于 $z_{tp}^\pm = \mp 1$ 处相撞。这两个 $u$ 值是 $\CB$ 中的
奇点。在每个奇点处，都有一个 BPS 粒子变成无质量的。当 $u$ 沿区间 $-\Lambda^2 < u < \Lambda^2$
变化时，两个转向点分开，绕 $z=0$ 朝相反方向运动，然后又重新会合。在该区间上任意一点
（例如 $u=0$）处的 BPS 谱由两个 BPS 超多重态组成。它们由连接两个转向点、但分别从
$z=0$ 处奇点的两侧经过的弦来描述。这两条弦共享两个端点，因此它们的电荷 $\gamma_{1,2}$
满足 $\langle \gamma_1, \gamma_2 \rangle = 2$。$\gamma_{1,2}$ 生成电荷格 $\hat\Gamma$。
$u$ 取在该区间内时典型的 WKB 三角剖分 $T_\WKB(\vartheta,u)$，以及它随 $\vartheta$
变化的变换，见图 57% TODO 合稿改 \ref{fig:Nf0plot}
与图 58% TODO 合稿改 \ref{fig:Nf0evolution-a}
所示。
\insfig{4.2in}{Nf0plot}{带 $N_f=0$ 的 $SU(2)$ 理论的一个典型 $T_\WKB$。为清楚起见，$z=\infty$ 已被映到有限距离。}
\insfig{6in}{Nf0evolution-a}{带 $N_f=0$ 的 $SU(2)$ 理论中 $T_\WKB(\vartheta, u)$ 的演化，$u$ 取在 $u=0$ 附近的强耦合区域。为清楚起见，$z=0$ 与 $z=\infty$ 已被映到有限距离。这里有两次 flip，对应于两个 BPS 态 $\gamma_1$、$\gamma_2$，且 $\langle \gamma_1, \gamma_2\rangle=2$。一条有限 WKB 曲线是红色虚线，另一条是蓝色虚线。红色（蓝色）竖直虚线标示出由于红色（蓝色）BPS 态的存在而引起的 flip。}

$\CB$ 中存在唯一一面椭球形的边缘稳定墙，它穿过那两个奇点。这面墙内部的区域
（我们刚才讨论的正是它）也称为「强耦合」区域，而外部的区域则是弱耦合区域。


注意，本理论中所有的 WKB 三角剖分都不过是两条内部边横穿其间的环带——这正是我们在
第 6.6.3 节% TODO 合稿改 \ref{sec:vector-jump}
中讨论极限三角剖分时所需要的前提情形。当 $u$ 位于强耦合区域时，随 $\vartheta$ 的演化
很简单，如图 58% TODO 合稿改 \ref{fig:Nf0evolution-a}
所示。另一方面，在弱耦合区域中我们会遇到一个新特征，即一个 BPS 矢量多重态。
对适当的 $\vartheta_c$，会出现单参数族的闭合 WKB 曲线，对应于电荷为 $\gamma_1+\gamma_2$
的单个矢量多重态，如图 59% TODO 合稿改 \ref{fig:Nf0plot-limit}
所示。于是，若 $u$ 位于弱耦合区域，则当 $\vartheta$ 的相位变化时，
$T_\WKB(\vartheta,u)$ 会经历无穷多次初等 flip，对应于电荷形如 $(n+1)\gamma_1 + n \gamma_2$
的 BPS 超多重态谱，其 BPS 射线在 $Z_{\gamma_1 + \gamma_2}$ 的相位处聚集。在矢量多重态的
临界相位处，$T_\WKB$ 作一次 \emph{juggle} 变换。随后又是一列无穷多次 flip，对应于电荷形如
$n \gamma_1 + (n+1) \gamma_2$ 的 BPS 超多重态谱。因此我们总共遇到的是该理论标准的弱耦合
BPS 谱：一个电荷为 $\gamma_1+\gamma_2$ 的 W 玻色子，以及一座无穷高的双荷子之塔。

\insfig{4.5in}{Nf0plot-limit}{带 $N_f=0$ 的 $SU(2)$ 理论的极限 WKB 三角剖分。为清楚起见，$z=\infty$ 已被映到有限距离，并画出了几条额外的闭合 WKB 曲线。}

比较这面墙两侧的谱，我们就得到电荷满足 $\langle \gamma_1, \gamma_2\rangle=2$ 时重要的
穿墙公式：
\begin{equation}\label{eq:wcf-prod2}
\CK_{\gamma_1} \CK_{\gamma_2}
= \CK_{\gamma_2} \CK_{\gamma_1+2 \gamma_2}\CK_{2\gamma_1+3\gamma_2} \cdots  \CK_{\gamma_1+\gamma_2}^{-2} \cdots \CK_{3\gamma_1+2\gamma_2}\CK_{2\gamma_1+\gamma_2}\CK_{\gamma_1}.
\end{equation}

更确切地说，这是对这面墙的两个截面之一而言的穿墙公式。对另一个截面，相应的公式是
\begin{equation}
\CK_{\gamma_2}\CK_{-\gamma_1}
= \CK_{-\gamma_1} \CK_{-2\gamma_1+ \gamma_2}\CK_{-3\gamma_1+2\gamma_2} \cdots  \CK_{-\gamma_1+\gamma_2}^{-2} \cdots \CK_{-2\gamma_1+3\gamma_2}\CK_{-\gamma_1+2\gamma_2}\CK_{\gamma_2}.
\end{equation}
同样，只要考虑到奇点周围预期的单值变换，右侧出现的这两个谱是彼此相容的：
\begin{equation}
M_1 = (\gamma_1, \gamma_2) \to (\gamma_1, \gamma_2-2 \gamma_1),
\end{equation}
\begin{equation}
M_2 = (\gamma_1, \gamma_2) \to (\gamma_1+2 \gamma_2, \gamma_2).
\end{equation}


\subsection{$N_f=1$}\label{sec:nf1}

如果在 $SU(2)$ 规范理论中加上一个味，二次微分就变为
\begin{equation}
\lambda^2 = \left( \frac{\Lambda^2}{z^3} + \frac{3u}{z^2}+ \frac{2
\Lambda m}{z} + \Lambda^2 \right) (dz)^2.
\end{equation}
同样，这也可以容易地用第 3 节% TODO 合稿改 \ref{sec:physics}
的方法推出。现在 $z=\infty$ 附近的展开显示出一个带两条 WKB 射线的不规则奇点。
质量参数 $m$ 与该奇点处的留数重合：
\begin{equation}
\lambda \sim \left( \Lambda + \frac{m}{z} + \cdots \right) dz.
\end{equation}

共有三个转向点。若取 $m = 0$，则 $\CB$ 中有三个奇点，位于 $u^3 = \frac{1}{4}
\Lambda^6$ 处。在每个奇点处，都有一个 BPS 态变成无质量的。在 $u=0$ 附近、与这三个奇点
相交的「强耦合」区域中，BPS 谱恰好由这三个 BPS 态组成。它们由连接相邻转向点对、
并绕 $z=0$ 顺时针行进的弦实现。它们的电荷 $\gamma_{1,2,3}$ 构成电荷格 $\hat \Gamma$
的一组基，满足 $\langle \gamma_1, \gamma_2\rangle=\langle \gamma_2, \gamma_3\rangle=\langle
\gamma_3, \gamma_1\rangle=1$。此外，$\sum \gamma_i$ 可以绕过 $z=0$，因而对应于一个
纯味荷。
\insfig{4.5in}{Nf1plot}{带 $N_f=1$ 的 $SU(2)$ 理论的 $T_\WKB$。为清楚起见，$z=0$ 与 $z=\infty$ 已被映到有限距离。}
\insfig{5.5in}{Nf1evolution-a}{带 $N_f=1$ 的 $SU(2)$ 理论中，$T_\WKB(\vartheta, u)$ 随 $\vartheta$ 的演化，$u$ 取在 $u=0$ 附近的强耦合区域。为清楚起见，$z=\infty$ 已被映到有限距离。这里有三次 flip，对应于电荷为 $\gamma_1$、$\gamma_2$、$-\gamma_3$ 的三个 BPS 态。}

当 $m=0$ 时，$Z_{-\gamma_3} = Z_{\gamma_1 + \gamma_2}$，且存在唯一的一面边缘稳定墙，
三个 BPS 态的中心荷在该墙上对齐。这面墙穿过那三个奇点。

我们现在把考查这些例子时一直采用的逻辑颠倒过来。至此，我们一直是根据 $T_\WKB$
对 $\vartheta$ 的依赖来推断 BPS 谱，然后再写出穿过一面墙时出现的那些辛同胚乘积之间的
恒等式。现在反过来，我们要用这些恒等式来确定 BPS 谱必定是什么。

例如，考虑这面墙的一个截面，其上 $Z_{\gamma_1}$、$Z_{\gamma_2}$、$Z_{-\gamma_3}$ 对齐。
在强耦合区域中，当 $\vartheta$ 越过相应的那些相位时，三角剖分经历三次 flip，它们合起来
给出辛同胚 $\CK_{\gamma_1} \CK_{-\gamma_3} \CK_{\gamma_2}$。在另一侧，同一个辛同胚必须
分解为相位按相反次序排列的因子。

为找到这样的分解，我们可以拿已经知道的恒等式来摆弄。把 $\CK_{\gamma_1}$ 搬到右边
毫无问题，只需用两次五边形恒等式
\begin{equation}
\CK_{\gamma_1} \CK_{-\gamma_3} \CK_{\gamma_2} = \CK_{-\gamma_3} \CK_{\gamma_1-\gamma_3}
\CK_{\gamma_2} \CK_{\gamma_1+\gamma_2}\CK_{\gamma_1}.
\end{equation}
然而，各因子仍未按恰当的次序排好：我们需要把 $\CK_{\gamma_2}$ 搬到左边，跨过
$\CK_{-\gamma_3} \CK_{\gamma_1-\gamma_3}$。我们有 $\langle \gamma_1-\gamma_3, \gamma_2 \rangle=2$，
因此需要用我们在 $N_f=0$ 理论中遇到过的恒等式 \eqref{eq:wcf-prod2}（作一次基变换），得到
\begin{multline}
\CK_{-\gamma_3} \CK_{\gamma_1-\gamma_3}\CK_{\gamma_2} \CK_{\gamma_1+\gamma_2}\CK_{\gamma_1}
= \\ \CK_{-\gamma_3} \CK_{\gamma_2} \CK_{2 \gamma_2 +\gamma_1-\gamma_3} \CK_{3 \gamma_2 +2\gamma_1-2\gamma_3}\cdots
\CK_{\gamma_1-\gamma_3+ \gamma_2}^{-2}\cdots \\ \cdots \CK_{2\gamma_2 +3\gamma_1-3\gamma_3}\CK_{\gamma_2 +2\gamma_1-2\gamma_3}\CK_{\gamma_1-\gamma_3} \CK_{\gamma_1+\gamma_2}\CK_{\gamma_1}.
\end{multline}
我们差不多做完了。因子 $\CK_{-\gamma_3}$ 与 $\CK_{\gamma_1+\gamma_2}$ 对应的电荷满足
$Z_{-\gamma_3} = Z_{\gamma_1 + \gamma_2}$，因此在正确重排后的乘积中，它们应当与电荷为
$\gamma_1 - \gamma_3 + \gamma_2$ 的矢量多重态相邻。要把它们搬到那里，我们需要让每一个
都穿过无穷多个其他因子；每一步都用五边形恒等式，最终得到
\begin{multline}
\CK_{\gamma_1} \CK_{-\gamma_3} \CK_{\gamma_2} = \\
\CK_{\gamma_2}\CK_{\gamma_2-\gamma_3} \CK_{2 \gamma_2 +\gamma_1-\gamma_3}\CK_{2 \gamma_2 +\gamma_1-2\gamma_3} \CK_{3 \gamma_2 +2\gamma_1-2\gamma_3}\CK_{3 \gamma_2 +2\gamma_1-3\gamma_3}\cdots
\CK_{-\gamma_3} \CK_{\gamma_1-\gamma_3+ \gamma_2}^{-2} \CK_{\gamma_1+\gamma_2}\cdots \\ \cdots \CK_{2\gamma_2 +3\gamma_1-3\gamma_3}\CK_{2\gamma_2 +3\gamma_1-2\gamma_3}
\CK_{2\gamma_2 +3\gamma_1-2\gamma_3}\CK_{\gamma_2 +2\gamma_1-2\gamma_3}\CK_{2\gamma_1+\gamma_2-\gamma_3}\CK_{\gamma_1-\gamma_3}\CK_{\gamma_1}.
\end{multline}
% 原文此式右侧乘积中 \CK_{2\gamma_2+3\gamma_1-2\gamma_3} 出现了两次，系原文如此，照录。
这确实对应于墙外弱耦合区域中的 BPS 谱。注意，电荷为 $-\gamma_3$ 与 $\gamma_1 + \gamma_2$
的两个「电」超多重态，是带 $N_f=1$ 的 $SU(2)$ 理论的拉氏量中的基本物质粒子，它们具有
相同的规范荷，为 W 玻色子的一半。

当我们在弱耦合区域中改变 $\vartheta$ 时，$T_\WKB$ 经历一列无穷多次 flip，然后作一个
类似于我们在第 6.6.3 节% TODO 合稿改 \ref{sec:vector-jump}
中描述过、并在 $N_f = 0$ 理论中遇到过的「\emph{juggle}」的变换。不过它并不完全相同；
特别是在极限三角剖分中，环带的其中一条边界上有 \ti{两} 个顶点（来自进入 $z=0$ 处奇点的
两条 WKB 射线）。如图 62% TODO 合稿改 \ref{fig:Nf1plot-limit}
所示。这是非一般的情形，它在这里出现是因为我们取了特殊值 $m=0$。这与下述事实有关：跨越
临界 $\vartheta$ 的坐标变换必须是
$\CK_{-\gamma_3} \CK_{\gamma_1-\gamma_3+ \gamma_2}^{-2} \CK_{\gamma_1+\gamma_2}$，
对应于一个矢量多重态加两个超多重态；而在通常情形下，临界 $\vartheta$ 处只有矢量多重态。

\insfig{5in}{Nf1plot-limit}{带 $N_f=1$ 且 $m=0$ 的 $SU(2)$ 理论的极限 $T_\WKB$。为清楚起见，$z=0$ 与 $z=\infty$ 已被映到有限距离，并画出额外的一条闭合 WKB 曲线。}

\insfig{6in}{Nf1plot-limit-pert}{带 $N_f=1$ 且 $m \neq 0$ 的 $SU(2)$ 理论的极限 $T_\WKB$。为清楚起见，$z=0$ 与 $z=\infty$ 已被映到有限距离，并画出额外的一条闭合 WKB 曲线。}

一旦我们把 $m$ 稍微偏离零，就回到一个更常规的格局。当我们从强耦合区域移向大 $u$ 时，
我们所描述的那一列初等操作（或与之非常接近的操作）发生在一列边缘稳定墙上。在很大但
有限的 $u$ 处，与 $\CK_{-\gamma_3}$、$\CK_{\gamma_1+\gamma_2}$ 相联系的那些超多重态的
中心荷，与和 $\CK_{\gamma_1-\gamma_3+ \gamma_2}^{-2}$ 相联系的中心荷并不完全对齐。
因此在重排乘积时，这两个因子 $\CK_{-\gamma_3}$ 与 $\CK_{\gamma_1+\gamma_2}$ 只应移动
有限多步。于是，围绕 $\CK_{\gamma_1-\gamma_3+ \gamma_2}^{-2}$ 的无穷乘积的「核心」，
对应于如第 6.6.3 节% TODO 合稿改 \ref{sec:vector-jump}
中那样的标准 flip 与 \emph{juggle} 序列。极限 WKB 三角剖分如图 63% TODO 合稿改 \ref{fig:Nf1plot-limit-pert}
所示；此时环带的两条边界上都只有一个顶点。

有意思的是把 $m$ 取到 $\frac{3}{2} \Lambda$ 附近，此时 $\CB$ 中有两个奇点合并。我们预期
第 9.4.4 节% TODO 合稿改 \ref{sec:N3}
中的 $N=3$ 理论会在两个奇点附近的标度区域中涌现。确实，看看标度区域中 WKB 叶层化的
形状（如图 64% TODO 合稿改 \ref{fig:Nf1plotN3}
所示），就会看到五边形出现在三个转向点正在汇聚的区域附近，这很有启发性。
\insfig{4.5in}{Nf1plotN3}{带 $N_f=1$ 的 $SU(2)$ 理论的 WKB 三角剖分，其中 $m$ 被调到接近
标度极限，在该极限下第 9.4.4 节% TODO 合稿改 \ref{sec:N3}
的 $N=3$ 理论涌现。}


\subsection{$N_f=2$，第一种实现}

利用第 3 节% TODO 合稿改 \ref{sec:physics}
的膜构造，有两条途径给 $SU(2)$ 规范理论加上第二个味。相应地，存在两个不同的 Hitchin
系统，而它们在 $\CM$ 上的度量预期是相同的。在本小节中我们处理对称的那一种可能性，
此时二次微分为
\begin{equation}
\lambda^2 = \left( \frac{\Lambda^2}{z^4} + \frac{2 \Lambda m_1}{z^3} + \frac{4u}{z^2} +  \frac{2 \Lambda m_2}{z}+ \Lambda^2 \right) dz^2
\end{equation}
于是 $z = 0$ 或 $z = \infty$ 附近的展开都显示出一个带两条 WKB 射线的不规则奇点。
两个质量参数 $m_{1,2}$ 是两个极点处的留数。

共有四个转向点，以及 $\CB$ 中的四个奇点。读者大概会猜到：再一次，$\CB$ 中将存在一个
强耦合区域，其中 BPS 谱只由那四个在四个奇点处无质量的 BPS 粒子组成。

若取 $m_{1,2} = 0$，就会出现一个惊人的简化。带 $N_f = 2$ 理论的 Seiberg--Witten
微分，乃至整个 Hitchin 系统，都可以通过简单的坐标变换 $z^2 \to \hat z$ 映到
$N_f = 0$ 理论的相应对象上。这使我们能够借用在 $N_f = 0$ 情形下的许多讨论。
$N_f = 0$ 理论的每个转向点都映为 $N_f = 2$ 理论的两个转向点；$N_f=0$ 理论的每个
BPS 超多重态都映为 $N_f=2$ 理论中的一对 BPS 超多重态。如图 65% TODO 合稿改 \ref{fig:Nf2plot-a}
所示。这一对超多重态具有相同的规范荷，但它们的味荷不同：两条相应路径之差可以形变为绕 $0$
与绕无穷各走一圈的路径之和。我们把其中一对的电荷记作 $\gamma_1^{1,2}$，另一对记作
$\gamma_2^{1,2}$，并有 $\langle \gamma_1^i, \gamma_2^j \rangle = 1$。此外
$\langle \gamma_1^i, \gamma_1^j\rangle =\langle
\gamma_2^i, \gamma_2^j\rangle =0$。若打开质量，则
$Z_{\gamma_1^1 - \gamma_1^2} = m_1+m_2$，而 $Z_{\gamma_2^1 -
\gamma_2^2} = m_1 - m_2$。

\insfig{3.2in}{Nf2plot-a}{带 $N_f=2$ 的 $SU(2)$ 理论第一种实现、$m_1 = m_2 = 0$ 时一个典型的 WKB 三角剖分。为清楚起见，$z=0$ 与 $z=\infty$ 已被映到有限距离。}
\insfig{3.2in}{Nf2plot-a-limit}{出现矢量多重态的极限 WKB 三角剖分 $T_\WKB(\vartheta_c)$，带 $N_f=2$ 的 $SU(2)$ 理论第一种实现、$m_1 = m_2 = 0$。为清楚起见，$z=\infty$ 已被映到有限距离；图中画出了代表该矢量多重态的单参数族中的两条闭合 WKB 曲线。}

按照这一推理，我们可以猜出相应穿墙公式的形状——它应当是 $N_f = 0$ 的公式
\eqref{eq:wcf-prod2} 的某种加倍，形如
\begin{equation}
\CK_{\gamma^1_1} \CK_{\gamma^2_1} \CK_{\gamma^1_2} \CK_{\gamma^2_2}
= \CK_{\gamma_2}\CK_{\gamma_2} \CK_{\gamma_1+2 \gamma_2}\CK_{\gamma_1+2 \gamma_2} \cdots\ \text{???}\ \cdots \CK_{2\gamma_1+\gamma_2}\CK_{2\gamma_1+\gamma_2}\CK_{\gamma_1}\CK_{\gamma_1}.
\end{equation}
\zy{原文该式右侧此处印作「???」，系作者留下的占位符，未补全；译文照录。}
在右侧，每个 $\gamma_{1,2}$ 究竟代表哪个循环 $\gamma_{1,2}^i$，先验地并不明显，
尽管可以通过考察 WKB 三角剖分中的 flip 来确定。另外，什么应当取代 \emph{juggle}
变换也不明显，因为环绕环状区域的每个圆上都有两个顶点——如图 66% TODO 合稿改 \ref{fig:Nf2plot-a-limit}
所示。当然，我们可以干脆拿已经知道的那些穿墙公式来摆弄，从而完全确定右侧。我们先把
$\CK_{\gamma_1^2}$ 搬到右边，得到
\begin{equation}
\CK_{\gamma_1^1}\CK_{\gamma_1^2} \CK_{\gamma_2^1}\CK_{\gamma_2^2} =
\CK_{\gamma_1^1} \CK_{\gamma_2^1}\CK_{\gamma_2^1+\gamma_1^2}\CK_{\gamma_2^2}\CK_{\gamma_2^2+\gamma_1^2} \CK_{\gamma_1^2},
\end{equation}
% TODO 原文此句作「and then do the same for $\CK_{\gamma_1^2}$」，而后一式实际是对
% $\CK_{\gamma_1^1}$ 作同样的操作；疑原文笔误，照译并在此标注。
然后再对 $\CK_{\gamma_1^2}$ 作同样的操作，得到
\begin{multline}
\CK_{\gamma_1^1} \CK_{\gamma_2^1}\CK_{\gamma_2^1+\gamma_1^2}\CK_{\gamma_2^2}\CK_{\gamma_2^2+\gamma_1^2} \CK_{\gamma_1^2}
= \\ \CK_{\gamma_2^1} \CK_{\gamma_2^1+\gamma_1^1} \CK_{\gamma_2^1+\gamma_1^2} \CK_{\gamma_2^1+\gamma_1^2+\gamma_1^1} \CK_{\gamma_2^2} \CK_{\gamma_2^2+\gamma_1^1}
\CK_{\gamma_2^2+\gamma_1^2}\CK_{\gamma_2^2+\gamma_1^2+\gamma_1^1}  \CK_{\gamma_1^2}  \CK_{\gamma_1^1}.
\end{multline}
我们还需要把 $\CK_{\gamma_2^2}$ 搬到左边，但它会跨过 $\CK_{\gamma_2^1+\gamma_1^2+\gamma_1^1}$，
从而产生一个如同我们在 $N_f = 0$ 例子中遇到过的无穷乘积：
\begin{multline}
\CK_{\gamma_2^1} \CK_{\gamma_2^1+\gamma_1^1} \CK_{\gamma_2^1+\gamma_1^2} \CK_{\gamma_2^1+\gamma_1^2+\gamma_1^1} \CK_{\gamma_2^2} \CK_{\gamma_2^2+\gamma_1^1}
\CK_{\gamma_2^2+\gamma_1^2}\CK_{\gamma_2^2+\gamma_1^2+\gamma_1^1}  \CK_{\gamma_1^2}  \CK_{\gamma_1^1}
= \\ \CK_{\gamma_2^1} \CK_{\gamma_2^1+\gamma_1^1} \CK_{\gamma_2^1+\gamma_1^2} \CK_{\gamma_2^2}\CK_{\gamma_2^1+2\gamma_2^2+\gamma_1^2+\gamma_1^1}\cdots
\CK_{\gamma_2^1+\gamma_2^2+\gamma_1^2+\gamma_1^1}^{-2}\cdots \\ \cdots
 \CK_{2\gamma_2^1+\gamma_2^2+2\gamma_1^2+2\gamma_1^1}\CK_{\gamma_2^1+\gamma_1^2+\gamma_1^1} \CK_{\gamma_2^2+\gamma_1^1}
\CK_{\gamma_2^2+\gamma_1^2}\CK_{\gamma_2^2+\gamma_1^2+\gamma_1^1}  \CK_{\gamma_1^2}  \CK_{\gamma_1^1}.
\end{multline}
因子
$\CK_{\gamma_2^1+\gamma_1^1}\CK_{\gamma_2^1+\gamma_1^2}$ 与
$\CK_{\gamma_2^2+\gamma_1^2}\CK_{\gamma_2^2+\gamma_1^1}$
对应于与矢量多重态 $\CK_{\gamma_2^1+\gamma_2^2+\gamma_1^2+\gamma_1^1}^{-2}$
相位相同的态。因此必须让它们穿过那个无穷乘积，其间用到一列五边形恒等式，最终得到
\begin{multline}
\CK_{\gamma_1^1}\CK_{\gamma_1^2} \CK_{\gamma_2^1}\CK_{\gamma_2^2}
= \\ \CK_{\gamma_2^1}  \CK_{\gamma_2^2}\CK_{\gamma_1^1+\gamma_2^2+\gamma_2^1} \CK_{\gamma_1^2+\gamma_2^2+\gamma_2^1}\CK_{\gamma_1^2+\gamma_1^1+\gamma_2^2+2\gamma_2^1}
\CK_{\gamma_1^2+\gamma_1^1+\gamma_2^1+2\gamma_2^2}\cdots \\ \cdots\CK_{\gamma_2^1+\gamma_1^1} \CK_{\gamma_2^1+\gamma_1^2}
\CK_{\gamma_2^1+\gamma_2^2+\gamma_1^2+\gamma_1^1}^{-2}\CK_{\gamma_2^2+\gamma_1^1}
\CK_{\gamma_2^2+\gamma_1^2} \cdots \\ \cdots \CK_{\gamma_1^2+2\gamma_1^1+\gamma_2^2+\gamma_2^1}
\CK_{2\gamma_1^2+\gamma_1^1+\gamma_2^2+\gamma_2^1}\CK_{\gamma_1^2+\gamma_1^1+\gamma_2^1} \CK_{\gamma_1^2+\gamma_1^1+\gamma_2^2}  \CK_{\gamma_1^2}  \CK_{\gamma_1^1}.
\end{multline}
与 $N_f = 1$ 的例子一样，若打开 $m_{1,2} \neq 0$，那唯一的一面墙就会碎成若干面墙。
在这些参数上玩一玩，可以得到一些有趣的结果。例如，若取 $m_1 = m_2 = 2 \Lambda$，
我们就得到一个漂亮的例子，
\begin{equation}
\lambda^2 = \left( \frac{\Lambda^2(z+1)^4}{z^4} + \frac{16 \tilde u}{z^2} \right) dz^2,
\end{equation}
其中 $\CB$ 的四个奇点里有三个在 $\tilde u = 0$ 处合并，并且我们能够恢复 $N=4$
Argyres--Douglas 标度极限。强耦合区域中四个粒子之一变得非常重，而另外三个与 $N=4$
Argyres--Douglas 理论的基本谱重合。

% ---------------------------------------------------------------------
% §10.4 插叙（原文 7479–7509 行；BRIEF 行段未覆盖，合稿时补译）
% ---------------------------------------------------------------------
\subsection{插叙：非阿贝尔味对称与 WKB 三角剖分} \label{subsec:intermission-flavor}

我们在第 3.2.8 节已经注意到% TODO 合稿改 \ref{sec:flavor}，当正则奇点 $\CP$ 处的质量参数趋于零、非阿贝尔 $SU(2)$ 味对称被恢复时，WKB
流线会有什么样的行为。考虑该极限下整片 WKB 叶状结构的行为是有益的。当 $m \to 0$ 时，一个转向点
$\CT$ 向 $\CP$ 移动。在一般的 $\vartheta$ 处，从 $\CT$ 出发的三条分离 WKB 曲线中有一条流向
$\CP$。另外两条朝相反方向绕过 $\CP$，然后彼此靠得很近地流走，最终落到某个别的奇点上。
（于是 $T_\WKB$ 是退化的，只有一条边从 $\CP$ 伸出。）

随着 $\vartheta$ 变化，我们偶尔会遇到相位非常接近的一对 BPS 态，此时这两条分界线变成有限
WKB 曲线并落到另一个转向点 $\CT'$ 上。这两个 BPS 态的电荷相差（两倍的）一个绕奇点的圈。
换言之，这些有限 WKB 曲线代表一对电荷为 $\gamma \pm \gamma_+$ 的 BPS 粒子二重态，其中
$\gamma_+$ 是与奇点 $\CP$ 相联系的味荷。

对处于 $Z_{\gamma \pm \gamma_+}$ 的相位之间那个狭窄窗口里的 $\vartheta$，这两条分界线从转向点
$\CT'$ 的两侧经过，并走向不同的奇点。WKB 三角剖分的相应变换是直截了当的。第一个 BPS 态
引起与 $\CP$ 相联系的退化三角形的非退化边发生一次 flip。于是有两条边终止于 $\CP$。第二个
BPS 态引起终止于 $\CP$ 的那条旧边发生一次 flip，于是我们又得到退化三角形。
\insfig{4.5in}{flavor-evolution}{当 $SU(2)$ 破缺质量很小时，$T_\WKB$ 中与一个 $SU(2)$ 味双重态
BPS 超多重态相关的两次 flip。}
